% Discussion paragraphs for §7
% Three standalone paragraphs: (1) Low F1 practical significance,
% (2) CF severity argument, (3) F1 significance difference explanation.
% All quantitative claims traced to registry exp_ids.

%%% Paragraph 1: Low F1 Practical Significance %%%
% exp_ids: phase2_scierc_full_s1 (R-F1=8.36%), phase0_scierc_unconstrained (baseline violations)

\paragraph{Practical Significance of Zero-Shot F1.}
Our reported relation F1 of approximately 8\% reflects the inherent difficulty of zero-shot closed IE, where the model receives no task-specific training examples.
For reference, supervised state-of-the-art systems on SciERC achieve 40--50\% F1 with full training data~\citep{zhong2021frustratingly,ye2022packed}.
Constraint enforcement operates orthogonally to extraction quality: it restricts the output space to ontologically valid triples regardless of the underlying model's accuracy.
When applied to fine-tuned models that reach competitive F1 levels, the same constraints remain necessary---Phase~0 analysis confirms that even well-prompted models exhibit non-trivial violation rates (15.37\% DR on SciERC).
% exp_id: phase0_scierc_unconstrained
In our evaluation, F1 serves as a \emph{control variable} verifying that constraint enforcement does not degrade extraction quality, not as an optimization target.
The contribution of this work is the constraint enforcement method itself; the zero-shot setting isolates its effect from confounds introduced by task-specific fine-tuning.

%%% Paragraph 2: CF Severity Argument %%%
% exp_ids: phase2_scierc_full_s1 (CF=9.90%), phase2_scierc_grammar_only (CF=0.08%),
%          phase2_docred_full_s1 (CF=2.57%)

\paragraph{Severity of Cross-Field Violations.}
Domain/range (DR) violations produce structurally invalid triples---a triple asserting that a \emph{Task} entity fills the \emph{domain} slot of \textsc{Part-of}, which requires \emph{Method}, cannot be consumed by any downstream system without error.
DR violations are therefore hard constraints.
Cross-field consistency (CF) violations are qualitatively different: each individual triple remains type-valid with respect to its own relation's domain/range rules; only cross-triple consistency is broken when the same entity receives different types across relations.
The dominant conflict pattern---\emph{Generic}$\leftrightarrow$\emph{OtherScientificTerm} (60 of 245~affected entities on SciERC)---reflects genuine semantic ambiguity in the type system rather than a systematic method failure.
% exp_id: phase2_scierc_full_s1
Ablation evidence cleanly attributes CF to the semantic layer: grammar-only decoding yields CF\,=\,0.08\%, matching the unconstrained baseline, while the full system reaches 9.90\% on SciERC and 2.57\% on DocRED.
% exp_ids: phase2_scierc_grammar_only, phase2_docred_full_s1
This arises from per-relation independent enforcement, where each relation's constraint mask is computed without access to type decisions made for other relations.
A natural extension is a global entity-level type consensus pass that resolves entity types before relation-specific enforcement.

%%% Paragraph 3: F1 Significance Difference Explanation %%%
% exp_ids: d012_bootstrap_ci (DocRED +0.50pp p<0.001, SciERC +0.13pp p=0.58),
%          phase0_docred_unconstrained (RV=29.66%), phase0_scierc_unconstrained (RV=1.39%)

\paragraph{Dataset-Dependent F1 Impact.}
The divergence in F1 significance between DocRED ($+$0.50\,pp, $p < 0.001$, 95\% CI [$+$0.23, $+$0.81]) and SciERC ($+$0.13\,pp, $p = 0.58$, 95\% CI [$-$0.24, $+$0.58]) traces to a single mechanism: grammar-layer elimination of out-of-vocabulary (OOV) relation names.
% exp_id: d012_bootstrap_ci
Under the all-valid-set evaluation protocol, a predicted triple whose relation name does not appear in the ontology is counted as a false positive.
The DocRED baseline exhibits a relation-vocabulary (RV) violation rate of 29.66\%---5{,}189 predictions with OOV relation names---all of which become false positives under this protocol.
% exp_id: phase0_docred_unconstrained
The grammar layer eliminates these OOV predictions, removing a substantial source of false positives and yielding a statistically significant precision gain.
On SciERC, the baseline RV rate is only 1.39\% (24~OOV predictions), too few to produce a detectable F1 effect.
% exp_id: phase0_scierc_unconstrained
The F1 improvement on DocRED is thus a direct, mechanistically transparent consequence of vocabulary enforcement at scale, not an artifact of constraint-induced distribution shift.
